% ============================================================================
% kapitel/emergenz.tex
% ============================================================================

\section{Die Emergenz der Physik aus Q}

Aus Q emergiert die gesamte Physik.

\section{Die DBT (De-Broglie-Bohm-Theorie)}

Die DBT ist die Quantentheorie.

Ihr Kern ist das Bohm-Potential:

\[
Q = -\frac{\hbar^2}{2m} \frac{\nabla^2 \sqrt{\rho}}{\sqrt{\rho}}
\]

Das ist Q in der DBT.

Die DBT ist Q.

\section{Die WG (Weber-Gravitation)}

Die WG ist die Gravitation.

Ihre Kraft ist:

\[
F_{\text{WG}} = -\frac{GMm}{r^2} \left[1 - \frac{\dot{r}^2}{c^2} + \frac{1}{2}\frac{r\ddot{r}}{c^2}\right]
\]

Aus Q folgt:

\[
F_{\text{WG}} = -\nabla Q
\]

Die WG ist der Gradient von Q.

\section{Die WED (Weber-Elektrodynamik)}

Die WED ist die Elektrodynamik.

Ihre Kraft ist:

\[
F_{\text{WED}} = \frac{q_1 q_2}{4\pi\epsilon_0 r^2} \left[1 - \frac{\dot{r}^2}{c^2} + 2\frac{r\ddot{r}}{c^2}\right]
\]

Aus Q folgt:

\[
F_{\text{WED}} = -\nabla Q
\]

Die WED ist der Gradient von Q.

\section{Die Einheit}

Die gesamte Physik ist Q.

\begin{itemize}
    \item DBT = Q
    \item WG = \( -\nabla Q \)
    \item WED = \( -\nabla Q \)
\end{itemize}

Q ist die Quelle.